\honest{Scope Insensitivity}{Eliezer Yudkowsky}{2007}

Three groups were asked how much they would pay to save 2,000, 20,000 or 200,000 migrating birds from drowning in oil ponds. They answered \$80, \$78 and \$88. I call this ``scope insensitivity'': the number of birds saved had little effect on what people would pay. \nb{The respondents were also told that the three numbers were ``much less than 1 percent,'' ``less than 1 percent'' and ``about 2 percent'' of the 8.5 million waterfowl in the flyway. Someone who cares about the population rather than each bird has less reason to pay a hundred times more for the largest number. The post does not mention this.}

Other surveys found the same. Toronto residents would pay little more to clean up all the polluted lakes in Ontario than to clean up those of one region. Residents of four western states would pay only 28\% more to protect all 57 wilderness areas in those states than to protect one. People picture ``a single exhausted bird, its feathers soaked in black oil, unable to escape,'' and the feeling this image produces sets what they pay. The image is the same whatever the number. No one can picture 2,000 birds, let alone 200,000, so the number is ``tossed out the window.'' \nb{The source, Kahneman, Ritov and Schkade, as Yudkowsky's own 2008 chapter on global risks quotes it, says the story ``probably evokes for many readers'' ``perhaps an image of an exhausted bird.'' The word ``single'' is the post's, and ``probably'' and ``perhaps'' are gone. The paper itself could not be checked here.} For what I call ``the usual finding,'' I give no source. It is that exponential increases in scope produce linear increases in payment. Each extra zero seems to add a small feeling to the image's feeling instead of multiplying it. \nb{The birds are the one study in the post with three levels, and their figures show no steady increase.} Only at the end of the paragraph do I call all this a hypothesis, ``valuation by prototype.''

A second hypothesis is that people buy a ``warm glow,'' whose price has nothing to do with the number of birds. I do not say which hypothesis the evidence favors.

Human lives are no different. When the risk from chlorinated drinking water was made about 600 times larger, what people would pay rose from \$3.78 to \$15.23. I call a fourfold rise insensitivity, and do not say what rise would have been right. Baron and Greene found no effect at all when the number of lives saved changed by a factor of ten.

My source for the claim that ``our perception of human deaths follows Weber's Law'' is a paper on ``psychophysical numbing.'' On that law, we notice a difference in deaths only when it is a fixed fraction of the total. So the same number of lives seems to matter less when the total is larger. \nb{The paper mentions Weber's law only as background.} A program to save 4,500 refugees won more support in a camp of 11,000 than in a camp of 250,000. \nb{In 2021, three replications of that study found the opposite preference; two replications of a related study in the same paper found the original effect.} And people said that a cure had to save far more lives to deserve funding when the disease killed more people each year. \nb{When the treatments were described by the lives they would save, 16\% of respondents answered this way, against over 67\% in the version the post reports.}

My moral: an effective altruist must think ``with the part of your brain that processes those unexciting inky zeroes on paper.'' \nb{Every study here records what people say they would pay or how they judge proposed programs. Kahneman, Ritov and Schkade, cited in the post, argue that such dollar answers are expressions of attitudes, not economic preferences. The post applies them to real giving.}
